\begin{textsample}{POS Dim 3 – ro – Score 25.00 – ro/ro\_inf\_27.txt}  \label{ex:f3_pos_135}
\textbf{Sugestões} de intencionalidades educativas para as práticas cotidianas de docentes na garantia dos objetivos de aprendizagens e desenvolvimento que serão construídos pelas \textbf{crianças} nas experiências e vivências, à medida que docentes oportunizem situações em que as \textbf{crianças} : - Compreendem o corpo em movimento como instrumento expressivo e de construção de novos conhecimentos de si, do outro e do universo. - Participem de momentos de preparar o ambiente interno e externo da unidade de Educação \textbf{Infantil} de modo que se sintam instigadas a explorá -los, por exemplo, transformando uma mesa em cabana, ou criando um túnel com caixas grandes de papelão. - Vivenciem com suas diferentes singularidades ( deficiências, transtornos globais de desenvolvimento e altas habilidades/superdotação ) sendo inclusas nos diversos momentos cotidianos. - Interpretem os \textbf{gestos} das outras \textbf{crianças} em interação comunicativa e/ou expressiva, verbalizando para elas sua compreensão do significado desses \textbf{gestos}. - Explorarem repetidamente os materiais, o espaço e o seu corpo de diferentes formas, com crescente domínio dos movimentos em danças e em representações teatrais. - Participem de jogos que envolvam orientar -se corporalmente — em frente, atrás, no alto, embaixo, dentro, fora —, em resposta a comandos dados por outras \textbf{crianças} ou pelo professor. - Recriem jogos acrescentando um desafio motor a um jogo já existente ( como jogar \textbf{futebol} com uma bola \textbf{menor} ) ou um conteúdo simbólico a um jogo de regra ( por exemplo, transformar um jogo de pega-pega em “ pega-monstro ” ). - \textbf{Brinquem} de esconde-esconde, de jogar bola, de pique, de seguir o mestre, de lenço atrás, de caça ao tesouro, de estátua, de barra-manteiga, de cabra-cega, de pula-sela, de pião etc. - \textbf{Manipulem} e dão vida a objetos, \textbf{brinquedos}, bonecos e fantoches em jogos teatrais. - Andem como robôs, zumbis, gatinhos ou maria-mole, entre outras formas. - \textbf{Batem}, esfregam, sopram, chacoalham objetos em \textbf{brincadeiras} ou canções, percebendo os movimentos corporais que realizam. - Criem histórias e narrativas e as \textbf{dramatizam} com os \textbf{colegas}, apropriando -se de diferentes gestualidades expressivas. - Dancem ao \textbf{som} de músicas de diferentes gêneros, \textbf{imitando}, criando e coordenando seus movimentos com os dos companheiros, usando diferentes materiais ( lenços, bola, fitas, instrumentos etc. ), explorando o espaço ( em cima, embaixo, para a frente, para trás, à esquerda e à direita ) e as qualidades do movimento ( rápido ou lento, forte ou leve ) a partir de \textbf{estímulos} diversos ( proposições orais, demarcações no chão, \textbf{mobiliário}, divisórias no espaço etc. ). - Fruem, descrevem, avaliam e reproduzem apresentações de dança de diferentes gêneros e outras expressões da cultura corporal ( circo, esportes, \textbf{mímica}, teatro etc. ), feitas por \textbf{adultos} amadores e profissionais ou por outras \textbf{crianças}. - Participem de danças como bumba meu boi, frevo, baião, maracatu, catira e outras do patrimônio indígena, afro-brasileiro, nipônico, italiano, alemão, boliviano etc., reproduzindo os movimentos e cantos, compreendendo o significado das indumentárias e das \textbf{pinturas} corporais utilizadas. - Construam em grupo roteiros para encenações feitas a partir de histórias conhecidas, situações improvisadas ou criações coletivas. - Teatralizem histórias conhecidas para outras \textbf{crianças} e \textbf{adultos}, apresentando movimentos e expressões corporais adequados a suas composições. - Encenem histórias com bonecos, fantoches ou figuras de sombras destacando \textbf{gestos}, movimentos, voz, caráter dos \textbf{personagens} etc. - Confeccionem cenários e figurinos para os enredos a serem \textbf{dramatizados}. - Assistem a apresentações de teatro profissional e popular com fantoches, sombras ou atores e identificam os elementos básicos dos roteiros apresentados. - Comentem apresentações de teatro feitas por outras \textbf{crianças} em relação aos objetos, fantoches, sombras ou \textbf{personagens} do enredo. - Vivenciem coletivamente a produção criativa de \textbf{gestos} e movimentos. - Participem de diferentes manifestações culturais e \textbf{brincadeiras} tradicionais. - Explorarem o corpo, o espaço e as primeiras coreografias improvisadas, ampliando o repertório de dança. - Explorem as possibilidades e limites da \textbf{movimentação} do corpo nas diferentes ações do cotidiano.
Todavia, torna -se necessário a disponibilidade de recursos para fortalecer as interações e \textbf{brincadeira} das \textbf{crianças}.
É importante que as \textbf{instituições} de Educação \textbf{Infantil} tenham acervo de \textbf{brinquedos} e \textbf{livros} que forneçam temas e ideias para o faz de conta, e com materiais como panos, caixas, blocos, madeiras, bem como, dispor de maneira criativa os móveis, utensílios, iluminação, \textbf{tecidos}, caixas, \textbf{brinquedos} que poderão ser usados pelas \textbf{crianças} para criar uma variedade de ambientes, cenários e \textbf{brincadeiras}.

% matched lemmas: adulto, bater, brincadeira, brincar, brinquedo, colega, criança, dramatizar, estímulo, futebol, gesto, imitar, infantil, instituição, livro, manipular, mobiliário, movimentação, mímico, pequeno, personagem, pintura, som, sugestão, tecido
\end{textsample}
